% Part: first-order-logic
% Chapter: sequent-calculus
% Section: proof-theoretic-notions

\documentclass[../../../include/open-logic-section]{subfiles}

\begin{document}

\begin{editorial}
  ఈ విభాగం సీక్వెంట్ కలనానికి సంబంధించిన నిరూపణీయతా సంబంధం,
  అవైరుధ్యం నిర్వచనాలను ఒకచోట చేర్చుతుంది.
  \sourcecorrection{OLTESEQ-003}{ఇది సీక్వెంట్ కలనం అధ్యాయంలోని
  విభాగం అయినప్పటికీ, ఆంగ్ల సంపాదకీయ గమనిక పొరపాటున ఈ నిర్వచనాలు
  “సహజ నిగమనానికి” చెందినవని చెప్పింది. విభాగపు గుర్తింపులు,
  $\Log{LK}$ నిర్వచనాలు, అధ్యాయ సందర్భం ప్రకారం దానిని “సీక్వెంట్
  కలనానికి”గా సరిచేశాం.}
\end{editorial}

\iftag{FOL}
      {\olfileid{fol}{seq}{ptn}}
      {\olfileid{pl}{seq}{ptn}}

\olsection{నిరూపణ-సిద్ధాంత భావనలు}

\begin{explain}
కొన్ని ముఖ్యమైన అర్థపర భావనలను (చెల్లుబాటుతనం, అనుగమనం,
సంతృప్తిపరచదగినతనం) నిర్వచించినట్లే, ఇప్పుడు వాటికి సరిపోలే
\emph{నిరూపణ-సిద్ధాంత భావనలను} నిర్వచిస్తాం. వీటిని \tetoken{నిర్మాణాలలో}{!!{structure}s}
\tetoken{వాక్యాలు}{!!{sentence}s} సంతృప్తి చెందడాన్ని ఆశ్రయించి నిర్వచించం; కొన్ని
సీక్వెంట్ల \tetoken{వ్యుత్పాద్యత}{!!{derivability}} లేదా \tetoken{అవ్యుత్పాద్యతను}{!!{nonderivability}} ఆశ్రయించి
నిర్వచిస్తాం. ఈ రెండు రకాల భావనలు ఏకీభవిస్తాయనేది ఒక ముఖ్యమైన
ఆవిష్కరణ. అవి ఏకీభవిస్తాయనే విషయమే \emph{నిర్దుష్టత},
\emph{సంపూర్ణతా సిద్ధాంతాల} సారం.
\end{explain}

\begin{defn}[సిద్ధాంతాలు]
\tetoken{ఒక వాక్యం}{!!^a{sentence}}~$!A$కు $\quad \Sequent !A$ సీక్వెంట్ యొక్క ఒక
$\Log{LK}$ \tetoken{ఒక వ్యుత్పత్తి}{!!a{derivation}} ఉంటే, దానిని \emph{సిద్ధాంతం} అంటాం.
$!A$ సిద్ధాంతమైతే $\Proves !A$ అనీ, కాకపోతే $\Proves/ !A$ అనీ రాస్తాం.
\end{defn}

\begin{defn}[\tetoken{వ్యుత్పాద్యత}{!!^{derivability}}]
\tetoken{వాక్యాలు}{!!{sentence}s} సమితి~$\Gamma$ నుంచి \tetoken{ఒక వాక్యం}{!!^a{sentence}}~$!A$
\emph{\tetoken{వ్యుత్పాదించదగిన}{!!{derivable}}}, అంటే $\Gamma \Proves !A$, అవుతుంది అంటే:
$\Gamma_0 \subseteq \Gamma$ అనే ఒక పరిమిత ఉపసమితి, $\Gamma_0$లోని
\tetoken{వాక్యాలతో}{!!{sentence}s} కూడిన $\Gamma_0'$ అనే ఒక క్రమం ఉండి, $\Log{LK}$
$\Gamma_0' \Sequent !A$ను \tetoken{నిగమనం}{!!{derive}s} చేయాలి. $\Gamma$ నుంచి $!A$
\tetoken{వ్యుత్పాదించదగిన}{!!{derivable}} కాకపోతే $\Gamma \Proves/ !A$ అని రాస్తాం.
\end{defn}

సంకోచనం, బలహీనీకరణ, మార్పిడి నియమాల వల్ల $\Gamma_0'$లోని
\tetoken{వాక్యాలు}{!!{sentence}s} క్రమమూ, వాటి సంఖ్యా ముఖ్యం కాదు: $\Gamma_0' \Sequent !A$
సీక్వెంట్ \tetoken{వ్యుత్పాదించదగిన}{!!{derivable}} అయితే, $\Gamma_0'$లో ఉన్న \tetoken{వాక్యాలే}{!!{sentence}s}
కలిగిన ఏ $\Gamma_0''$కైనా $\Gamma_0'' \Sequent !A$ కూడా
నిగమించదగినది. ఉదాహరణకు $\Gamma_0 = \{!B, !C\}$ అయితే,
$\Gamma_0' = \tuple{!B, !B, !C}$, $\Gamma_0'' =
\tuple{!C, !C, !B}$ రెండూ $\Gamma_0$లోని \tetoken{వాక్యాలే}{!!{sentence}s} కలిగిన
క్రమాలు. వీటిలో ఒకదాన్ని కలిగిన సీక్వెంట్ \tetoken{వ్యుత్పాదించదగిన}{!!{derivable}} అయితే
మరొకదాన్ని కలిగినదీ అలాగే అవుతుంది; ఉదాహరణకు:
\begin{prooftree}
  \AxiomC{}
  \Deduce$!B, !B, !C \fCenter !A$
  \RightLabel{\LeftR{\Contraction}}
  \UnaryInf$!B, !C \fCenter !A$
  \RightLabel{\LeftR{\Exchange}}
  \UnaryInf$!C, !B \fCenter !A$
  \RightLabel{\LeftR{\Weakening}}
  \UnaryInf$!C, !C, !B \fCenter !A$
\end{prooftree}
ఇకపై $\Gamma_0$ ఒక పరిమిత \tetoken{వాక్యాలు}{!!{sentence}s} సమితి అయితే,
$\Gamma_0 \Sequent !A$ అనేది $\Gamma_0$లోని \tetoken{వాక్యాలు}{!!{sentence}s} క్రమాన్ని
పూర్వాంగంగా కలిగిన ఏ సీక్వెంట్‌నైనా సూచిస్తుందని అంటాం; అవసరమైతే
సంకోచనాలు, మార్పిడులు, బలహీనీకరణలను మౌనంగా చేర్చినట్లు భావిస్తాం.

\begin{defn}[అవైరుధ్యం]
$\Gamma_0 \subseteq \Gamma$ అనే ఏదో పరిమిత ఉపసమితికి
$\Gamma_0 \Sequent \quad$ను $\Log{LK}$ \tetoken{నిగమనం}{!!{derive}s} చేస్తే, అప్పుడు
మాత్రమే వాక్యాల సమితి~$\Gamma$ \emph{వైరుధ్యభరితం}. $\Gamma$
వైరుధ్యభరితం కాకపోతే, అంటే ప్రతి పరిమిత $\Gamma_0 \subseteq \Gamma$కు
$\Log{LK}$ అనేది $\Gamma_0 \Sequent \quad$ను \tetoken{నిగమనం}{!!{derive}} చేయకపోతే,
దాన్ని \emph{అవైరుధ్యమైనది} అంటాం.
\end{defn}

\begin{prop}[స్వావర్తనత్వం]
\ollabel{prop:reflexivity}
$!A \in \Gamma$ అయితే $\Gamma \Proves !A$.
\end{prop}

\begin{proof}
ప్రారంభ సీక్వెంట్ $!A \Sequent !A$ అనేది \tetoken{వ్యుత్పాదించదగిన}{!!{derivable}}; అంతేకాక
$\{!A\} \subseteq \Gamma$.
\end{proof}

\begin{prop}[ఏకదిశత]
\ollabel{prop:monotonicity}
$\Gamma \subseteq \Delta$, $\Gamma \Proves !A$ అయితే
$\Delta \Proves !A$.
\end{prop}

\begin{proof}
$\Gamma \Proves !A$ అనుకుందాం. అంటే $\Gamma_0 \subseteq \Gamma$ అనే
ఒక పరిమిత సమితికి $\Gamma_0 \Sequent !A$ \tetoken{వ్యుత్పాదించదగిన}{!!{derivable}} అవుతుంది.
$\Gamma \subseteq \Delta$ కనుక $\Gamma_0$ అనేది $\Delta$కు కూడా
పరిమిత ఉపసమితి. కాబట్టి $\Gamma_0 \Sequent !A$ యొక్క \tetoken{వ్యుత్పత్తి}{!!{derivation}}
$\Delta \Proves !A$ అని కూడా చూపుతుంది.
\end{proof}

\begin{prop}[సంక్రామకత్వం]
\ollabel{prop:transitivity}
$\Gamma \Proves !A$, $\{!A\} \cup \Delta \Proves !B$ అయితే
$\Gamma \cup \Delta \Proves !B$.
\end{prop}

\begin{proof}
$\Gamma \Proves !A$ అయితే $\Gamma_0 \subseteq \Gamma$ అనే ఒక పరిమిత
సమితి, $\Gamma_0 \Sequent !A$ యొక్క \tetoken{ఒక వ్యుత్పత్తి}{!!a{derivation}}~$\pi_0$ ఉంటాయి.
$\{!A\} \cup \Delta \Proves !B$ అయితే ఏదో పరిమిత ఉపసమితి
$\Delta_0 \subseteq \Delta$కు $!A, \Delta_0 \Sequent !B$ యొక్క
\tetoken{ఒక వ్యుత్పత్తి}{!!a{derivation}}~$\pi_1$ ఉంటుంది. కింది \tetoken{వ్యుత్పత్తిని}{!!{derivation}} పరిశీలించండి:
\begin{prooftree}
  \AxiomC{}
  \RightLabel{$\pi_0$}
  \Deduce$\Gamma_0 \fCenter !A$
  \AxiomC{}
  \RightLabel{$\pi_1$}
  \Deduce$!A, \Delta_0 \fCenter !B$
  \RightLabel{\Cut}
  \BinaryInf$\Gamma_0, \Delta_0 \fCenter !B$
\end{prooftree}
$\Gamma_0 \cup \Delta_0 \subseteq \Gamma \cup \Delta$ కనుక ఇది
$\Gamma \cup \Delta \Proves !B$ అని చూపుతుంది.
\end{proof}

ప్రత్యేకంగా $\Gamma \Proves !A$, $!A \Proves !B$ అయితే
$\Gamma \Proves !B$ అని దీని అర్థమని గమనించండి. అలాగే
$!A_1, \dots, !A_n \Proves !B$ అయి, ప్రతి~$i$కి
$\Gamma \Proves !A_i$ అయితే $\Gamma \Proves !B$ అని కూడా వస్తుంది.

\begin{prop}
\ollabel{prop:incons}
ప్రతి వాక్యం~$!A$కు $\Gamma \Proves {!A}$ అయినప్పుడు, అప్పుడు మాత్రమే
$\Gamma$ వైరుధ్యభరితం.
\end{prop}

\begin{proof}
అభ్యాసం.
\end{proof}

\begin{prob}
\olref[fol][seq][ptn]{prop:incons}ను నిరూపించండి.
\end{prob}

\begin{prop}[సంహతత్వం]
  \ollabel{prop:proves-compact}
  \begin{enumerate}
  \item $\Gamma \Proves !A$ అయితే $\Gamma_0 \Proves !A$ అయ్యేలా
    $\Gamma_0 \subseteq \Gamma$ అనే ఒక పరిమిత ఉపసమితి ఉంటుంది.
  \item $\Gamma$ యొక్క ప్రతి పరిమిత ఉపసమితి అవైరుధ్యమైతే,
    $\Gamma$ అవైరుధ్యమైనది.
  \end{enumerate}
\end{prop}

\begin{proof}
  \begin{enumerate}
    \item $\Gamma \Proves !A$ అయితే $\Gamma_0 \subseteq \Gamma$ అనే
      ఒక పరిమిత ఉపసమితికి $\Gamma_0 \Sequent !A$ సీక్వెంట్ యొక్క
      \tetoken{ఒక వ్యుత్పత్తి}{!!a{derivation}} ఉంటుంది. కాబట్టి $\Gamma_0 \Proves !A$.
    \item $\Gamma$ వైరుధ్యభరితమైతే, $\Gamma_0 \subseteq \Gamma$ అనే
      ఒక పరిమిత ఉపసమితికి $\Gamma_0 \Sequent \quad$ను $\Log{LK}$
      \tetoken{నిగమనం}{!!{derive}s} చేస్తుంది. కానీ అప్పుడు $\Gamma_0$ అనేది
      వైరుధ్యభరితమైన, $\Gamma$ యొక్క పరిమిత ఉపసమితి.
  \end{enumerate}
\end{proof}

\end{document}
